\documentclass[border=5mm]{standalone}
% ==============================================================================
% SETUP.TEX - CONFIGURACIÓN GLOBAL DEL PROYECTO
% ==============================================================================
% Este archivo centraliza todos los paquetes y estilos.
% Debe incluirse en el preámbulo del libro principal y de los archivos standalone.
% \PassOptionsToPackage{gray}{xcolor}
% --- 1. CODIFICACIÓN E IDIOMA ---
% fix-cm habilita las fuentes Computer Modern en tamaños arbitrarios (por debajo
% de 5 pt, entre otros). Debe cargarse antes que fontenc.
\usepackage{fix-cm}
\usepackage[utf8]{inputenc}
\usepackage[T1]{fontenc}
\usepackage[spanish, es-tabla]{babel}  
\usepackage{csquotes} % Recomendado para babel y citas

\usepackage{import}
% mode=image: Usa el PDF pre-compilado, nunca intenta recompilar.
% Debes compilar cada figura manualmente antes de compilar el libro.
\usepackage[mode=image, subpreambles=false]{standalone}

% --- 2. MATEMÁTICAS Y CIENCIAS ---
\usepackage{amsmath, amssymb, amsfonts, mathtools}
\usepackage{enumitem}

% LaTeX trae \DeclareMathSizes{5}{5}{5}{5}, así que a 5 pt (\tiny) los subíndices
% salen del mismo tamaño que la base: en las etiquetas de figuras el M de
% $\omega_M$ queda desproporcionado. Se restablece la proporción habitual.
\DeclareMathSizes{5}{5}{3.7}{3}

% --- 3. DISEÑO Y GEOMETRÍA --- 
% Unificamos los márgenes para todo el documento
\usepackage[left=2.5cm,right=2.5cm,top=3cm,bottom=3cm]{geometry}
\makeatletter
\ifstandalone % Evita que geometry dañe el recorte de las figuras bajándolas 3cm
  \@ifclasswith{standalone}{crop=false}{
    % Es un capítulo/sección: Dejamos que geometry actúe.
  }{
    % Es una figura: Neutralizamos geometry para que no las recorte mal.
    \geometry{pass}
  }
\fi
\makeatother
\usepackage{parskip} % Espaciado entre párrafos sin sangría
\usepackage{float}   % Para usar [H] en figura

% Ajustes de flotantes para evitar espacios verticales exagerados
\renewcommand{\topfraction}{0.95}      % Permite que las figuras ocupen hasta el 95% superior
\renewcommand{\bottomfraction}{0.95}   % Permite que las figuras ocupen hasta el 95% inferior
\renewcommand{\textfraction}{0.05}     % Permite hojas con tan solo un 5% de texto
\renewcommand{\floatpagefraction}{0.8} % Requiere que una página de solo figuras esté 80% llena
\setcounter{topnumber}{4}
\setcounter{bottomnumber}{4}
\setcounter{totalnumber}{10}

% --- 4. GRÁFICOS Y TABLAS ---
\usepackage{graphicx}
\usepackage{booktabs} % Tablas profesionales
\usepackage{caption}  % Control de captions y \captionof
\usepackage{tikz}
\usetikzlibrary{arrows.meta, calc, angles, quotes, babel, backgrounds}
\usepackage{pgfplots}
\usepgfplotslibrary{groupplots}
\usepgfplotslibrary{external}
\usepgfplotslibrary{patchplots}
\pgfplotsset{compat=1.18}
\usepackage[american, straightvoltages]{circuitikz}

% --- 5. COLORES Y CAJAS ---

\usepackage[table]{xcolor}
\usepackage{tcolorbox}

% Definición de colores corporativos/estilo
\definecolor{utnblue}{RGB}{0, 51, 102}
\definecolor{codegreen}{rgb}{0,0.6,0}
\definecolor{codegray}{rgb}{0.5,0.5,0.5}
\definecolor{codepurple}{rgb}{0.58,0,0.82}
\definecolor{backcolour}{rgb}{0.96,0.96,0.96}

% --- 6. ENTORNOS PERSONALIZADOS (CAJAS) ---

% Caja "Resultado" (Azul Corporativo) — Definiciones, fórmulas clave, resultados importantes.
% Uso: \begin{resultado}[Fórmula de Euler] ... \end{resultado}
\newtcolorbox{resultado}[1][]{
    colback=utnblue!5!white,
    colframe=utnblue,
    title=\textbf{#1},
    fonttitle=\bfseries,
    sharp corners,
    boxrule=0.5mm
}

% Caja "Nota" (Gris) — Advertencias, aplicaciones, contexto secundario.
% Uso: \begin{nota}[Cuidado con la calculadora] ... \end{nota}
% Uso: \begin{nota} ... \end{nota}  → título por defecto "Nota"
\newtcolorbox{nota}[1][Nota]{
    colback=codegray!5!white,
    colframe=codegray,
    title=\textbf{#1},
    fonttitle=\bfseries,
    sharp corners,
    boxrule=0.5mm
}

% Caja inline para resaltar un valor y enlazarlo con su otra aparición en una tabla
% o en una cuenta: mismo color, mismo valor.
% Uso: \marcavalor{$4$}  o  \marcavalor[codegreen]{$4$}
\newtcbox{\marcavalor}[1][utnblue]{
    on line,
    boxsep=1pt, left=2pt, right=2pt, top=1pt, bottom=1pt,
    colback=#1!10!white,
    colframe=#1,
    boxrule=0.4pt,
    arc=2pt
}

% --- 7. CÓDIGO FUENTE (LISTINGS) ---
\usepackage{listings}

\lstdefinestyle{miestilo}{
    backgroundcolor=\color{backcolour},   
    commentstyle=\color{codegreen},
    keywordstyle=\color{magenta},
    numberstyle=\tiny\color{codegray},
    stringstyle=\color{codepurple},
    basicstyle=\ttfamily\footnotesize,
    breakatwhitespace=false,         
    breaklines=true,                 
    captionpos=b,                    
    keepspaces=true,                 
    numbers=left,                    
    numbersep=5pt,                  
    showspaces=false,                
    showstringspaces=false,
    showtabs=false,                  
    tabsize=2,
    frame=single,                    
    rulecolor=\color{codegray}
}

\lstset{style=miestilo}
% Soporte para tildes y ñ en bloques de código
\lstset{literate={ñ}{{\~n}}1 {á}{{\'a}}1 {é}{{\'e}}1 {í}{{\'i}}1 {ó}{{\'o}}1 {ú}{{\'u}}1}

% Operadores matemáticos personalizados

% Vector en sentido discreto (lista finita de coordenadas): negrita + flecha.
% Uso: \vect{v}, \vect{e}_k, \vect{0}.
\newcommand{\vect}[1]{\vec{\mathbf{#1}}}

% Fecha de versión y comando para capítulos standalone
\providecommand{\verdate}{\the\year.\ifnum\month<10 0\fi\the\month.\ifnum\day<10 0\fi\the\day}
\newcommand{\chapterversion}{%
    \vspace{-1.5cm}\begin{flushright}\small\textit{\color{codegray}Versión~\verdate}\end{flushright}\vspace{0.5cm}%
}

% --- 8. HIPERVÍNCULOS (DEBE SER EL ÚLTIMO PAQUETE) ---
\usepackage[colorlinks=true, linkcolor=utnblue, urlcolor=utnblue, citecolor=utnblue]{hyperref}

% --- 9. REFERENCIAS CRUZADAS ENTRE CAPÍTULOS STANDALONE ---
% Cuando se compila un capítulo standalone, xr-hyper lee los .aux de los
% otros capítulos (si existen) para resolver \ref{} a labels externos.
% En la compilación del libro completo este bloque no se activa.
\ifstandalone
  \usepackage{xr-hyper}
  \makeatletter
  \newcommand{\xrchap}[1]{%
    \edef\xrc@self{\detokenize{Capitulo#1}}%
    \edef\xrc@job{\jobname}%
    \ifx\xrc@self\xrc@job\else
      \IfFileExists{../#1capitulo/Capitulo#1.aux}%
        {\externaldocument{../#1capitulo/Capitulo#1}}{}%
    \fi
  }
  \makeatother
  \xrchap{01}\xrchap{02}\xrchap{03}\xrchap{04}
  \xrchap{05}\xrchap{06}\xrchap{07}\xrchap{08}
  \xrchap{09}\xrchap{10}\xrchap{11}\xrchap{12}
  \xrchap{13}
\fi
% \PassOptionsToPackage{gray}{xcolor}

\begin{document}
\begin{tikzpicture}[>=Latex, font=\small]

% Señal: x(t) = exp(-t/3)*sin(pi*t/2)
%
% Pulso base: p_Δ(t) = 1/Δ para 0 ≤ t < Δ  →  área = 1
% Aproximación de la señal: \hat{x}(t) = \sum x(k\Delta) p_\Delta(t-k\Delta) \Delta
% Cada pulso rectangular en la suma tiene:
%   altura = x(k\Delta)
%   ancho  = \Delta
%   área   = x(k\Delta) \Delta
%
% Al reducir \Delta, la suma aproxima la integral de convolución
%
% Valores altura = x(k\Delta) para cada panel:
%
%  Δ=1.00:  t=0 → 0.000    t=1 → 0.717    t=2 → 0.000    t=3 → -0.368
%
%  Δ=0.50:  t=0.0→0.000  t=0.5→0.599  t=1.0→0.717  t=1.5→0.429
%           t=2.0→0.000  t=2.5→-0.307 t=3.0→-0.368 t=3.5→-0.220
%
%  Δ=0.25:  t=0.00→0.000  t=0.25→0.352  t=0.50→0.599  t=0.75→0.720
%           t=1.00→0.717  t=1.25→0.609  t=1.50→0.429  t=1.75→0.214
%           t=2.00→0.000  t=2.25→-0.181 t=2.50→-0.307 t=2.75→-0.369
%           t=3.00→-0.368 t=3.25→-0.313 t=3.50→-0.220 t=3.75→-0.110

\pgfplotsset{
  approx/.style={
    width=5.5cm, height=4.5cm,
    xmin=-0.2, xmax=4.5,
    ymin=-0.65, ymax=1.35,
    axis lines=center,
    xlabel={$t$}, xlabel style={anchor=south west, inner sep=1pt},
    xtick={0,1,2,3,4},
    x tick label style={font=\scriptsize},
    ytick={-1,0,1,2,3},
    y tick label style={font=\scriptsize},
    tick style={thin},
    title style={font=\small, yshift=0ex},
    clip=false,
  },
  pulso/.style={
    utnblue, fill=utnblue!15, fill opacity=0.7,
    ybar interval, line width=0.4pt,
  },
  destacado/.style={
    utnblue, fill=utnblue!45, fill opacity=0.9,
    ybar interval, line width=0.7pt,
  },
  curva/.style={
    red, thick, densely dotted, domain=0:4, samples=200,
  },
}

% ============================================================
% Panel (a): Δ = 1   →   altura = x(k)
% ============================================================
\begin{axis}[approx, name=ax1,
    title={(a) $\Delta = 1$},
]

\addplot[pulso] coordinates {
    (0,0)(1,0.717)(2,0)(3,-0.368)(4,0)
};
% pulso destacado en t=[1,2]: área = x(1)*1 ≈ 0.717
\addplot[destacado] coordinates {(1,0.717)(2,0)};
% curva de referencia (punteada)
\addplot[curva] {exp(-x/3)*sin(deg(pi*x/2))};
% anotación ancho Δτ sobre el pulso destacado
\draw[|-|, red, thin] (axis cs:1,-0.45) --
    node[below, font=\scriptsize, inner sep=1pt]{$\Delta$} (axis cs:2,-0.45);
% anotación área del pulso destacado ubicada arriba
\node[font=\scriptsize, utnblue, anchor=south, inner sep=2pt] (area1) at (axis cs:1.5, 1.05)
    {área $= x(1)\Delta$};
\draw[->, utnblue, densely dashed, thin, shorten >=2pt, shorten <=1pt] 
    (area1.south) -- (axis cs:1.5, 0.717);

\end{axis}

% ============================================================
% Panel (b): Δ = 0.5   →   altura = x(kΔ)
% ============================================================
\begin{axis}[approx, name=ax2,
    at={(ax1.outer east)}, anchor=outer west, xshift=2mm,
    title={(b) $\Delta = 0{,}5$},
]

\addplot[pulso] coordinates {
    (0.0, 0.000)(0.5, 0.599)(1.0, 0.717)(1.5, 0.429)
    (2.0, 0.000)(2.5,-0.307)(3.0,-0.368)(3.5,-0.220)(4.0,0)
};
% pulso destacado en t=[1,1.5]: área x(1)*0.5 ≈ 0.358
\addplot[destacado] coordinates {(1.0,0.717)(1.5,0)};
\addplot[curva] {exp(-x/3)*sin(deg(pi*x/2))};
\draw[|-|, red, thin] (axis cs:1,-0.45) --
    node[below, font=\scriptsize, inner sep=1pt]{$\Delta$} (axis cs:1.5,-0.45);
\node[font=\scriptsize, utnblue, anchor=south, inner sep=2pt] (area2) at (axis cs:1.25, 1.05)
    {área $= x(1)\Delta$};
\draw[->, utnblue, densely dashed, thin, shorten >=2pt, shorten <=1pt] 
    (area2.south) -- (axis cs:1.25, 0.717);

\end{axis}

% ============================================================
% Panel (c): Δ = 0.25   →   altura = x(kΔ)
% ============================================================
\begin{axis}[approx, name=ax3,
    at={(ax2.outer east)}, anchor=outer west, xshift=2mm,
    title={(c) $\Delta = 0{,}25$},
]

\addplot[pulso] coordinates {
    (0.00, 0.000)(0.25, 0.352)(0.50, 0.599)(0.75, 0.720)
    (1.00, 0.717)(1.25, 0.609)(1.50, 0.429)(1.75, 0.214)
    (2.00, 0.000)(2.25,-0.181)(2.50,-0.307)(2.75,-0.369)
    (3.00,-0.368)(3.25,-0.313)(3.50,-0.220)(3.75,-0.110)
    (4.00,0)
};
% pulso destacado en t=[1,1.25]: área x(1)*0.25 ≈ 0.179
\addplot[destacado] coordinates {(1.00,0.717)(1.25,0)};
\addplot[curva] {exp(-x/3)*sin(deg(pi*x/2))};
\draw[|-|, red, thin] (axis cs:1,-0.45) --
    node[below, font=\scriptsize, inner sep=1pt]{$\Delta$} (axis cs:1.25,-0.45);
\node[font=\scriptsize, utnblue, anchor=south, inner sep=2pt] (area3) at (axis cs:1.125, 1.05)
    {área $= x(1)\Delta$};
\draw[->, utnblue, densely dashed, thin, shorten >=2pt, shorten <=1pt] 
    (area3.south) -- (axis cs:1.125, 0.717);

\end{axis}

\end{tikzpicture}
\end{document}
